\scanpage{7}
\mspage{6}
\unsure{``} \uline{Sous-famille} \inserted{(ou sous-ens)} \uline{admissible} d'un ens admissible $\unsure{\mathscr{L}|\mathcal{L}}\subset\unsure{\mathcal{P}|\mathscr{P}}(L)$ : par définition, c'est un $\Psi\subset\mathscr{L}$ qui est \uline{fermé} dans $\mathscr{L}$. Alors pour $A\in\Psi$, $\deleted{\unsure{\partial_{\mathscr{L}}A}}$ on a
\[
 \partial_{\Psi}A=\partial_{\mathscr{L}}A\quad\text{donc}\quad
 A^{\circ(\deleted{\mathscr{L}}\Psi)}=A^{\circ(\mathscr{L})},
\]
ce qui montre que $|\Psi|$ est une partie saturée de $|\mathscr{L}|$, pour la rel. d'équiv. $\deleted{\unsure{\mathscr{L}}}R_{\mathscr{L}}$ \unsure{\ensuremath{s=}|sur|\ensuremath{S=}} $|\mathscr{L}|$,
\deleted{La top. de $\Psi$ est induite}

En fait, avec les notations précédentes, on a
\[
 |\Psi|=A_\Psi
\]
\deleted{De \ensuremath{=}, il est évident que pour \unsure{ce|x}} Soit $x\in|\Psi|$, donc $\deleted{\unsure{A}}\exists A\in\Psi$ tel que $\unsure{a|x}\in A$, \deleted{$x\in|\Psi|\Rightarrow$} et donc $x\in A(x)\subset A$ (puisque $A(x)$ est le plus petit des éléments de $\mathscr{L}$ qui $\ni x$) donc $A(x)\in\Psi$ ($\Psi$ étant fermé). Donc $A(x)$ est aussi le plus petit élément \uline{de $\Psi$} qui contient $x$ :

\deleted{\unsure{Ain}} Ainsi, $\Psi$ est préadmissible, et de plus il est évident que il est admissible (les $A^{\circ\Psi}\ne\varnothing$, d'après $A^{\circ(\Psi)}=A^{\circ\mathscr{L}}$).

\uline{Prop.} Soit $\mathscr{L}\subset\mathcal{P}(L)$ ens de parties admissible, et soit $\Psi\subset\mathcal{P}(L)$. Conditions équivalentes

a) $\Psi\subset\mathscr{L}$, et $\Psi$ est une partie fermée de $\mathscr{L}$

b) $\Psi\subset\mathscr{L}$, $\Psi$ est admissible, \inserted{$|\Psi|$ est saturé dans $|\mathscr{L}|$ pour $\unsure{R_{\Psi}|R_{\varphi}|R_{\mathscr{L}}}$,} et la relation d'équivalence \inserted{$R_\Psi$} \unsure{sur|de} $|\Psi|\subset|\mathscr{L}|$ définie par $\Psi$ est celle induite par $R_{\mathscr{L}}$
